\documentclass[11pt,a4paper,twoside]{article}
\usepackage[utf8]{inputenc}
\usepackage[english]{babel}
\usepackage{actuarialangle, actuarialsymbol}


\begin{document}

Motivate, Define Choquet capacities and give examples

Bad events have small probability as they're usually tail events 
\( \lim_{x \to \infty} 1 - F_S(x) = \lim_{x \to \infty} \bar{F}_S(x) = 0\)
but an unproportional impact, thus they should be given more weight.
\( F \) could however stop being a probability measure under such modifications.
A framework which maintains a reasonable amount of previous properties
are Choquet properties

A Choquet capacity is a function \(\mu: \mathcal{F} \rightarrow [0, 1]\) such that

1. \(\mu(\emptyset) = 1 - \mu(\Omega) = 0\)
2. \(A, B \in \mathcal{F}, A \subset B \implies \mu(A) \le \mu(B)\)

(so its missing the \( \sigma \)-additivity of probability measures).
Examples are

Distorted probabilities: 
\(\mathbb{P}\) probability measure on \((\Omega, \mathcal{F})\), 
\(h: [0, 1] \rightarrow [0, 1]\) nondecreasing function with \(h(0) = 1 - h(1) = 0\)
\(\implies \mu := h \circ \mathbb{P}\) Choquet capacity

Upper/lower probabilities: 
\(\mathcal{P}\) set of probability measures on \((\Omega, \mathcal{F})\), define
\(\mu^*(A) = \sup_{\mathbb{P} \in \mathcal{P}} \mathbb{P}(A)\)
\(\mu_*(A) = \inf_{\mathbb{P} \in \mathcal{P}} \mathbb{P}(A)\),
both are Choquet capacities.
\(\mathcal{P}\) set of probability measures, \(\mu^*, \mu_*\) as before, 
\(\alpha \in [0, 1]\), then \(\mu := \alpha \mu^* + (1 - \alpha) \mu_*\) is a Choquet capacity.


Define Choque integrals and state some necessary lemma's

Lemma:
\(S \ge 0\) random variable on probability space \((\Omega, \mathcal{F}, \mathbb{P})\). Then
\(\mathbb{E}[S] = \int_{\mathbb{R}_{+}} \mathbb{P}(S > x) dx.\)

Proof.
Using Fubini's Theorem
\(\mathbb{E}[S] = \int_{\mathbb{R}_+} s \mathbb{P}_S(ds)\)
\(= \int_{\mathbb{R}_+} \int_{\mathbb{R}_+} \mathbf{1}_{(0, s)}(x) dx \mathbb{P}_S(ds)\)
\(= \int_{\mathbb{R}_+} \int_{\mathbb{R}_+} \mathbf{1}_{(0, s)}(x) \mathbb{P}_S(ds) dx \)
\(= \int_{\mathbb{R}_+} \int_{\mathbb{R}_+} \mathbf{1}_{(x, \infty)}(s) \mathbb{P}_S(ds) dx\)
\(= \int_{\mathbb{R}_+} \mathbb{P}(S > x) dx.\)

Lemma:
\( \mathbb{E}[S] = \int_{0}^\infty \mathbb{P}(S > x)dx + \int_{-\infty}^0 \mathbb{P}(S > x) - 1 dx \)

Let \( \mu : \mathcal{F} \to [0, 1] \) be a Choquet capacity and \( f \in \mathcal{L}^\infty \).
The Choquet integral of \( f \) with respect to \( \mu \) is:

\( \int f d\mu = \int_{0}^\infty \mu(f > x)dx + \int_{-\infty}^0 \mu(f > x) - 1 dx \)
\(x \ge |f|_\infty \implies \{f > x\} = \emptyset \implies \mu(f > x) = 0 \implies 0 \le \int_0^\infty \mu(f > x) dx \le |f|_\infty\)
\(x \le -|f|_\infty \implies \{f > x\} = \Omega \implies \mu(f > x) = 1 \implies 0 \ge \int_{-\infty}^0 (\mu(f > x) - 1) dx \ge -|f|_\infty\)

The Choquet integral w.r.t. a Choquet capacity \(\mu\) on \((\Omega, \mathcal{F})\) has the following properties:
1. Monotonicity: Suppose \(f, g \in L^\infty\) satisfy \(f(w) \le g(w), w \in \Omega\). Then
\(\int f d \mu \le \int g d \mu\)
2. Cash-additivity: For all \(f \in L^\infty\) and \(c \in \mathbb{R}\):
\(\int (f + c) d \mu = \int f d \mu + c.\)
3. Lipschitz continuity: For all \(f, g \in L^\infty\):
\(|\int f d \mu - \int g d \mu| \le \sup_{w \in \Omega} |f(w) - g(w)|\)
It does however not satisfy
\( \int fd\mu + \int g d\mu \neq \int (f + g) d\mu \)
however for comonotone functions, that is \( (f(s) - f(s'))(g(s) - g(s')) \geq 0 \) for all \( s, s' \in S \),
or that \( \mu \) is \( \sigma \)-additive, this becomes an equality.
\end{document}
